\documentclass[a4paper,10.5pt]{article}
\usepackage[T1]{fontenc}\usepackage[utf8]{inputenc}
\usepackage[top=2.2cm,bottom=2.2cm,left=2.35cm,right=2.35cm]{geometry}
\usepackage{hyperref}\usepackage{xcolor}\definecolor{accent}{HTML}{1F4E78}\hypersetup{colorlinks=true,urlcolor=accent}
\pagestyle{empty}\setlength{\parindent}{0pt}\setlength{\parskip}{7pt}
\begin{document}
\textbf{Prabhavit Gupta}\\
Venice, Italy \textbar{} +39 351 610 7807 \textbar{} prabhavitg@gmail.com\\[10pt]
Dear Blueocean Technologies Hiring Team,

I am writing to apply for the Forward Deployed Engineer role in Berlin. Over the past three years I have worked as the engineer who sits between a business problem and a working AI or automation solution, across eight European companies, and that is exactly the bridge this role asks for: sitting with the people who do the work, then building and owning the agent that fixes it.

At FlowSpace, I took over a live workspace marketplace and automated its data migration off a legacy stack onto Next.js and Laravel while the product stayed live, mapping the old structure to the new one and choosing the migration path per entity, either API endpoints or CSV exports. The platform grew from 2,000 to 4,500 plus users during that period. I also built and still own an AI listing bot that takes a workspace provider's brochure in any format and drafts a full listing with a human in the loop, more than 900 drafts so far, and I built support AI agents that answer customer queries from live platform data and escalate to a person on low confidence, the same production discipline of evaluation, structured output and graceful failure that owning agents in production requires.

At UltraStrategy, I built the CV parsing pipeline behind the platform. A user uploads a CV, the pipeline extracts the data with Affinda and scores how relevant the candidate is for specific role types, the result reaches the frontend, and the pipeline then generates a standardised Ultra CV using the Claude API. That pipeline was built with n8n and Eden AI. I also spent time on Get2Germany, integrating Google APIs and building automation tools for the German market, so I am comfortable being close to a local, domain-specific audience rather than a generic user base.

I like being close to the actual work rather than the abstraction of it. I am used to picking up context directly from the people doing the job, deciding what to build without a handed-down spec, and shipping something that survives contact with real use rather than winning an argument about architecture.

I hold Italian EU Long Term Residence Permit status. I am not currently authorised to work in Germany and would need the applicable employer side permit or sponsorship process. I am ready to relocate to Berlin and complete that process.

Regards,\\
Prabhavit Gupta
\end{document}
